\section{Smooth balls carrying physical volume}
\label{sec:physical-volume}

Let \(g_i\) be smooth unit-volume Yamabe minimizers on a fixed smooth
oriented homotopy four-sphere \(M\), with positive Yamabe constants
\(Y_i=R_{g_i}\) and \(E_+(g_i)\leq 2T-\delta\) for a fixed
\(\delta>0\), where \(T=4\pi^2\). Write \(\mu_i=dV_{g_i}\) and
\(\nu_i=|\operatorname{Rm}(g_i)|^2dV_{g_i}\), using the full tensor
norm for \(\operatorname{Rm}\). The curvature identities
\eqref{eq:budgets}--\eqref{eq:schouten} give
\[
 Y_0:=\sqrt{48\delta}\leq Y_i\leq S:=8\sqrt{6}\,\pi,
 \qquad \sup_i\nu_i(M)<\infty,
 \qquad \int_M\sigma_2(A_{g_i})\,dV_{g_i}\geq\delta/2.
\]
Here \(S\) is the round Yamabe constant and the upper bound is the
Yamabe comparison with the sphere \cite[Theorem~2]{Schoen1984Yamabe}.
With \(A_g=(\operatorname{Ric}_g-R_g g/6)/2\), the sign of the last
integral agrees with the sign in the doubled-Schouten convention of
Chang--Gursky--Yang. Their positive Yamabe and positive total
\(\sigma_2\) criterion gives a smooth metric of positive Ricci curvature
in each conformal class \cite[Theorem~A and Corollary~B]{ChangGurskyYang2002PositiveRicci}.

Normalize any such metric by a constant rescaling to obtain \(k_i\),
in which the least radius of a ball carrying half of \(\mu_i\) is one.
Denote a witnessing center by \(p_i\). For each fixed smooth metric
distance spheres have zero four-volume by polar integration, and
hence also zero \(\mu_i\)-mass. Ball mass is continuous in center and
radius by inclusions of nearby balls and continuity in radius.
Maximizing over the compact center set then proves attainment and
gives
\[
 \mu_i(B_1(p_i))=\tfrac12=\max_{x\in M}\mu_i(B_1(x)).
\]
In this section source balls and distances use \(k_i\), and on a
limit space they use its pointed Gromov--Hausdorff limit distance.

\begin{theorem}[Finite smooth volume witnesses]
\label{thm:volume-balls}
For the sequence above there is a subsequence on which the following
holds. For any \(0<\varepsilon<1\), \(\zeta>0\) and \(r_*>0\), there
are finitely many fixed radii \(r_\ell\in(0,r_*)\) and, for every
sufficiently large \(i\), pairwise disjoint smoothly embedded closed
standard four-balls \(\Omega_{i\ell}\) in \(M\), oriented by \(M\),
with
\begin{equation}\label{eq:physical-ball-witness}
 \Omega_{i\ell}\subset B_{r_\ell}(x_{i\ell}),\qquad
 \nu_i(B_{r_\ell}(x_{i\ell}))<\zeta,\qquad
 \mu_i\Bigl(\bigcup_\ell\Omega_{i\ell}\Bigr)>1-\varepsilon.
\end{equation}
The centers \(x_{i\ell}\) approach regular points of a single pointed
limit of the \(k_i\). The number of balls and the radii may depend on
the tolerances and on the subsequence. The theorem applies to arbitrary
sequences in the stated band, including ones with vanishing Bach defect.
\end{theorem}

\subsection{Capacity, mass bounds, and noncollapse}

For a Lipschitz cutoff \(0\leq\chi\leq1\) whose nonzero set has
\(\mu_i\)-mass at most \(1/2\), the physical Yamabe inequality and
H\"older's inequality give
\begin{equation}\label{eq:physical-small-support}
 Y_0(1-2^{-1/2})\,\mu_i(\{\chi=1\})^{1/2}
 \leq 6\int_M |d\chi|_{g_i}^2\,d\mu_i
 \leq 6\left(\int_M |d\chi|_{k_i}^4\,dV_{k_i}\right)^{1/2}
        \mu_i(U)^{1/2},
\end{equation}
where \(U\) contains the gradient support up to a null set. Indeed,
the scalar term in the Yamabe inequality is at most
\(Y_i2^{-1/2}(\int\chi^4d\mu_i)^{1/2}\). The last inequality uses
conformal invariance of the fourth-power gradient integral in
dimension four.

For \(0<r<R\) take a cutoff equal to one on \(B_r(x)\), affine in
the logarithm of distance on \(U=B_R(x)\setminus\overline B_r(x)\),
and zero outside \(B_R(x)\). By the Bishop radial area bound
\(2\pi^2t^3\), its fourth-power gradient integral is at most
\(2\pi^2/\log^3(R/r)\), as is that of the reversed cutoff. Squaring
\eqref{eq:physical-small-support} bounds the plateau mass by
\(C_0\mu_i(U)/\log^3(R/r)\), with
\(C_0=72\pi^2/[Y_0^2(1-2^{-1/2})^2]\). Use the inward cutoff when
\(R\leq1\) and the reversed cutoff when \(x=p_i\) and \(1\leq r\).
Each step of a fixed ratio \(e^L>1\) then contracts the inner mass,
or the outer tail mass, by the factor \(C_0/(C_0+L^3)<1\).
Iteration and monotonicity yield constants \(C,\alpha>0\), independent
of \(i\), such that
\begin{equation}\label{eq:physical-power}
 \sup_x\mu_i(B_r(x))\leq C r^\alpha\quad(0<r\leq1),\qquad
 \mu_i(M\setminus B_R(p_i))\leq C R^{-\alpha}\quad(R\geq1).
\end{equation}

Relative Bishop--Gromov comparison gives a packing bound even before
a uniform lower bound for \(k_i\)-volume. For an \(a\)-separated set
in \(B_D(p_i)\), the disjoint \(a/2\)-balls lie in \(B_{D+a/2}(p_i)\).
The ball of radius \(2D+a/2\) about each center contains this common
enclosing ball, so each disjoint ball has volume at least
\((1+4D/a)^{-4}\) times the volume of the enclosing ball. Thus a
maximal such net has at most \((1+4D/a)^4\) centers. With \(a=1/4\)
and \(D=1\), some covering \(1/4\)-ball has physical mass at least a
fixed \(m>0\). A linear distance cutoff supported in the concentric
\(1/2\)-ball has \(k_i\)-gradient at most \(4\), support mass at most
one half, and fourth-power gradient integral at most
\(4^4\operatorname{Vol}_{k_i}(B_2(p_i))\). Equation
\eqref{eq:physical-small-support} gives a uniform positive lower
bound for this last ball volume. For \(x\in B_D(p_i)\), the ball
\(B_{D+2}(x)\) contains \(B_2(p_i)\). Applying relative comparison
there and the Euclidean upper bound gives, for every fixed \(D\),
\begin{equation}\label{eq:physical-noncollapse}
 c_D r^4\leq \operatorname{Vol}_{k_i}(B_r(x))\leq\omega_4r^4,
 \qquad x\in B_D(p_i),\quad 0<r\leq1,\quad \omega_4=\pi^2/2.
\end{equation}

\subsection{The conformal factor on bounded windows}

Write \(g_i=w_i^2k_i\). With \(\Delta=\operatorname{div}\nabla\),
the scalar equation and \(R_{k_i}>0\) give
\(-\Delta_{k_i}w_i\leq V_iw_i\), where \(V_i=(Y_i/6)w_i^2\).
On \(B_r(x)\) with \(x\in B_D(p_i)\) and \(r\leq1\), the
local Sobolev inequality for nonnegative Ricci curvature
\cite[Theorem~3.1]{SaloffCoste1992Uniformly} has a prefactor
bounded by a dimensional constant times
\(r^2\operatorname{Vol}_{k_i}(B_r(x))^{-1/2}\). By
\eqref{eq:physical-noncollapse}, it gives
\[
 \|z\|_4^2\leq C_D\int_{B_r(x)} (|dz|^2+r^{-2}z^2)\,dV_{k_i},
 \qquad z\in H^1_0(B_r(x)).
\]
All norms in this argument use \(k_i\). Testing the scalar
subsolution with \(\chi^2w_i^{p-1}\) for \(p\geq2\), where
\(\chi\) is supported in the ball, gives the gradient term
\(4(p-1)p^{-2}\int\chi^2|d(w_i^{p/2})|^2dV_{k_i}\). Absorb the
mixed cutoff term by Young's inequality. The Sobolev inequality for
\(z=\chi w_i^{p/2}\) then gives
\begin{equation}\label{eq:physical-moser}
 \|z\|_4^2\leq C_Dp^2\left[
 \int(|d\chi|^2+r^{-2}\chi^2)w_i^p\,dV_{k_i}
 +\int V_i z^2\,dV_{k_i}\right].
\end{equation}
At \(p=4\) the potential term is at most
\((S/6)\mu_i(B_r(x))^{1/2}\|z\|_4^2\). Choose a fixed
\(r=r(D)>0\) small enough to absorb it by \eqref{eq:physical-power}.
Using a cutoff equal to one on \(B_{r/2}(x)\) and
\(\int_M w_i^4\,dV_{k_i}=1\) gives a uniform \(L^8\) bound for
\(w_i\) on \(B_{r/2}(x)\), hence an \(L^4\) bound for \(V_i\) there.

For subsequent cutoffs in \(B_{r/2}(x)\), interpolation bounds
\(\int V_i z^2\,dV_{k_i}\leq C_{D,r}\|z\|_2\|z\|_4\).
Absorbing the \(L^4\) factor in \eqref{eq:physical-moser} costs
at most \(C_{D,r}p^4\|z\|_2^2\). For exponents \(p_j=8\,2^j\)
and radii decreasing geometrically from \(r/2\) to \(r/4\),
the cutoff gradients are bounded by \(C2^j/r\), so the norm
multipliers are at most \([C_{D,r}p_j^4 4^j]^{1/p_j}\). Their
product is finite. Varying \(x\) in the window proves
\begin{equation}\label{eq:physical-density-bound}
 \sup_{B_D(p_i)}w_i\leq C_D',\qquad
 \mu_i\leq (C_D')^4 dV_{k_i}\quad\text{on } B_D(p_i).
\end{equation}

\subsection{Regular points and transfer to the source}

We now prove Theorem~\ref{thm:volume-balls}. The packing estimate
gives a subsequential pointed Gromov--Hausdorff limit \((X,d,p)\).
By \eqref{eq:physical-noncollapse} and the Ricci lower bound,
Riemannian volumes converge locally weakly to Euclidean-normalized
\(m=\mathcal H^4\) on \(X\)
\cite[Theorem~5.9]{CheegerColding1997StructureI}. In a common pointed
realization take local weak limits \(\mu_i\rightharpoonup\mu\) and
\(\nu_i\rightharpoonup\nu\) after a further subsequence. The tail
bound gives \(\mu(X)=1\), and \eqref{eq:physical-density-bound}
gives \(\mu=fm\) with \(f\in L^\infty_{\mathrm{loc}}(X)\).
Cheeger--Colding regularity gives Euclidean four-dimensional tangent
cones at \(m\)-almost every point
\cite[Theorems~2.1 and 5.11]{CheegerColding1997StructureI}.
At such a point \(m(B_r(x))/(\omega_4r^4)\to1\). To see the density
claim, approximate a sequence of blow-ups with rescaled sources
chosen by a diagonal argument using fixed-radius volume convergence.
Their unit volumes stay uniformly positive by
\eqref{eq:physical-noncollapse}. Noncollapsed volume convergence
to the Euclidean tangent then gives the ratio limit.

The measure \(m\) is locally doubling by relative volume comparison.
Differentiation of \(\mu\) with respect to \(m\) and the fact that
\(\mu\) gives zero mass to the countable atoms of \(\nu\) show that
at \(\mu\)-almost every \(x\) the following hold, for each fixed \(A>1\).
If \(x_i\to x\) in the pointed realization, then
\begin{equation}\label{eq:physical-ordered-limits}
 \lim_{r\downarrow0}\lim_{i\to\infty}
      \frac{\mu_i(B_r(x_i))}{\omega_4 r^4}=f(x)\in(0,\infty),
 \qquad
 \lim_{r\downarrow0}\lim_{i\to\infty}\nu_i(B_{Ar}(x_i))=0.
\end{equation}
The radii in these iterated limits are continuity radii for the
corresponding limit measures. In particular one first transfers
fixed radii by weak convergence and then shrinks in \(X\).

Choose a bounded window with \(\mu\)-mass greater than
\(1-\varepsilon/4\), and \(0<\tau<1/6\) with
\((1-3\tau)^4(1-\varepsilon/2)>1-\varepsilon\).
The Cheeger--Colding almost-maximal-volume embedding theorem says
that, for sufficiently small radius and volume deficit, a ball
\(B_s(x_i)\) of a smooth manifold with the stated Ricci lower bound
contains a smoothly embedded open model ball whose image contains
\(B_{(1-\tau)s}(x_i)\)
\cite[Theorems~A.1.8--A.1.9]{CheegerColding1997StructureI}.
These thresholds are fixed once \(\tau\) is fixed. Each smooth
source qualifies as a constant noncollapsed sequence for the theorem.

At the regular positive-density points of the window that are
outside the atoms of \(\nu\), take a fine family of closed limit
balls of radii \(r<r_*\) with \(\nu(\overline B_r(x))<\zeta/2\).
By the density statements, choose arbitrarily small radii for
which the \(m\)-volume ratio at \(s=(1-\tau)r\) meets the embedding
threshold with strict margin, and for which
\(\mu(B_{(1-3\tau)r}(x))/\mu(B_r(x))\) is as close to
\((1-3\tau)^4\) as needed in the strict inequality above.
Choose \(r,s,(1-3\tau)r\) at joint continuity radii.
The Vitali covering theorem for locally doubling \(m\) selects
disjoint balls covering almost all of these centers with respect to
\(m\) and hence to \(\mu\). A finite subfamily then carries mass
greater than \(1-\varepsilon/2\). Its smaller
\((1-3\tau)r_\ell\)-balls still carry mass greater than
\(1-\varepsilon\).

Hold this finite family and its radii fixed and transfer the centers
to the sources. The \(s_\ell=(1-\tau)r_\ell\) balls have a positive
separation margin inherited from the disjoint outer limit balls,
and their volumes meet the smooth embedding criterion for large
\(i\). Each image contains the closed concentric ball of radius
\((1-3\tau)r_\ell\) compactly, since
\((1-3\tau)<(1-\tau)^2\). The preimage of that closed ball is
compact in the open model ball. A smaller closed concentric model
ball containing this preimage in its interior has image
\(\Omega_{i\ell}\). A reflection of the parameter ball fixes the
orientation if needed. Weak convergence on the fixed inner and
outer balls now proves \eqref{eq:physical-ball-witness} for all
sufficiently large indices on this subsequence.

In the Hilbert--Schmidt operator convention
\(E_+(g_i;U)\leq\nu_i(U)/4\). If \(\zeta<4\eta\), the interior of
each sphere in \eqref{eq:physical-ball-witness} has chiral mass
below \(\eta\), the prescribed lower mass on each side of an
\(\eta\)-balanced cut. The bounds are consistent with a complement
carrying nearly all the chiral energy in small physical volume.
The density estimate concerns \(\mu_i\), and the selection
explicitly avoids curvature atoms. To use a cut for a minimizing
sequence with \(E_+(g_i)\to I^+(M)\), one also needs a compatible
marking of its two-sided buffer in a gauge with the first-order
estimates needed for both positive-Yamabe caps. The smooth ball
markings supply spherical faces and radial room in \(k_i\).

\subsection{Selecting the two measures together}

On this same pointed subsequence, also take the local weak limit
\(\lambda\) of \(\lambda_i=|W^+_{g_i}|^2_{\rm HS}dV_{g_i}\).
Let \(\mathcal R\) be the Euclidean regular set of the limit of
the positive-Ricci comparison metrics, and put
\(G_\zeta=\{x\in\mathcal R:\nu(\{x\})<\zeta\}\).
In particular \(\mu(G_\zeta)=1\). Write
\(e_\infty=\lim_i e_i-\lambda(X)\) for the mass lost at infinity.

\begin{proposition}\label{prop:joint-vol}
At fixed \(\varepsilon,\zeta,r_*>0\) the finite family in
Theorem~\ref{thm:volume-balls} may also be chosen with
\[
 \limsup_i\lambda_i(M\setminus\textstyle\bigcup_j\Omega_{ij})
 \le e_\infty+\lambda(X\setminus G_\zeta)+\varepsilon .
\]
This includes chiral mass singular to physical volume on \(G_\zeta\).
\end{proposition}
\begin{proof}
Use the finite measure \(\sigma=\mu+\lambda\). Cover all but
\(\varepsilon/4\) of its mass on \(G_\zeta\) by a bounded window.
There its packing number at radius \(r\) is at most \(C_Dr^{-4}\).
For dyadic \(r\), the set of centers where
\(\sigma(B_r(x))<r^5\) is covered by at most \(C_Dr^{-4}\) such
balls, using a separated net in the set. Summability and
monotonicity give \(\sigma(B_r(x))\ge 2^{-5}r^5\) at all small
radii at almost every center for \(\sigma\).

For any fixed small \(\eta>0\), choose \(\theta=1-3\tau\) close
enough to one. At those centers there are arbitrarily small
continuity radii with
\[
 \sigma(B_r)\le(1+\eta)\sigma(B_{\theta r}),\qquad
 \sigma(B_{5r})\le C_\theta\sigma(B_r).
\]
Indeed set \(r_j=r_0\theta^j\), avoiding the countably many bad
radii for the measures and their needed dilates. The nondecreasing
numbers \(a_j=-\log\sigma(B_{r_j})\) satisfy
\(a_j\le C_x+5j\log(1/\theta)\). The upper proportion of jumps
\(a_{j+1}-a_j>\log(1+\eta)\) is at most
\(5\log(1/\theta)/\log(1+\eta)<1/4\). Choose \(k\) with
\(\theta^{-k}\ge5\). Summing the \(k\)-step jumps counts every
one-step jump at most \(k\) times, so the upper proportion with
\(a_j-a_{j-k}>\log C_\theta\) is less than \(1/4\) for large enough
\(C_\theta\). This proves the two simultaneous bounds.

Restrict this fine family to \(G_\zeta\) and impose the outer
curvature ceiling, the regular-point comparison-volume requirement,
and the maximum radius from the proof of Theorem~\ref{thm:volume-balls}.
A disjoint greedy selection by dyadic size has fivefold dilates
covering the centers. Centers outside the selected closed balls are
covered by the fivefold dilates of every remaining tail, since they
have arbitrarily small admissible balls missing any finite initial
union. The \(\sigma\)-mass of each tail of dilates is bounded by
\(C_\theta\) times the sum of the disjoint tail masses, hence tends
to zero. Keep a finite subfamily. Shrinking the selected balls to
\(\theta r\) loses at most \(\eta\sigma(X)\). Choose \(\eta\) and
the tail loss within the tolerance. At these fixed radii the same
marked embeddings contain the transferred inner balls and satisfy
the strict outer curvature ceiling. Their union captures all but
\(\varepsilon\) of \(\mu(G_\zeta)\) and of \(\lambda(G_\zeta)\)
in the source lower limits. Subtract from the total limiting chiral
mass to obtain the inequality.
\end{proof}
